% GENERATED FILE. Edit canonical lessons, then run npm run generate:books.
% canonical-source-hash: fnv1a64:ca5b33c053bf8b15
% canonical-lessons: IT-C442-salita, IT-C442-discesa, IT-C442-costeggiare, IT-C442-sbucare, IT-C442-aggirare

\chapter{Su e giù}
\label{ch:it-su-e-giu}
\hlchaptermodality{442}

\begin{chapteropening}
\textbf{By the end of this chapter:} \emph{I can describe a route: uphill, downhill, running alongside, coming out somewhere and going round.}

Complete the last lesson of chapter 442: I can describe a route: uphill, downhill, running alongside, coming out somewhere and going round.
\end{chapteropening}

\section[la salita]{la salita — the climb, the uphill slope}

\label{lesson:IT-C442-salita}



\begin{quote}
Before the new one: say the Italian for the avenue, then the Italian for the one-way street.
\end{quote}

\subsection*{You'll want to know: la salita}
\textbf{la salita} — \textquotedblleft{}the climb, the uphill slope\textquotedblright{}.

\begin{quote}
La salita fino al castello è dura.
\end{quote}

\begin{quote}
\emph{The climb up to the castle is hard.}
\end{quote}

\emph{In salita} — uphill.

\begin{cousinweb}
\emph{salita} is built on \emph{salire} 'to go up', which is inherited from Latin \emph{salīre} 'to jump' (EtymDB).
\end{cousinweb}

\begin{grammarlens}[title={What you've built}]
One more word for telling, explaining and describing at B1.
\end{grammarlens}

\subsection*{Your turn}
\begin{itemize}
\raggedright
  \item \emph{Say it:} \emph{la salita}
  \item \emph{Say it:} \emph{la salita}, once more
  \item \emph{Say it:} a sentence of your own with \emph{la salita}
  \item \emph{Recall:} say the Italian for the avenue, then the Italian for the one-way street, then say \emph{la salita} again
\end{itemize}

\begin{tcolorbox}[breakable,colback=teal!4,colframe=teal!35!black,title={Before you move on}]
What does la salita mean? (\textquotedblleft{}The climb, the uphill slope\textquotedblright{}.) Say it once more.
\end{tcolorbox}
\section[la discesa]{la discesa — the descent, the downhill slope}

\label{lesson:IT-C442-discesa}



\begin{quote}
Before the new one: say the Italian for the one-way street, then the Italian for the climb.
\end{quote}

\subsection*{You'll want to know: la discesa}
\textbf{la discesa} — \textquotedblleft{}the descent, the downhill slope\textquotedblright{}.

\begin{quote}
In discesa la bici va da sola.
\end{quote}

\begin{quote}
\emph{Downhill the bike goes by itself.}
\end{quote}

\emph{In discesa} — downhill; on the slopes, \emph{la discesa} is the downhill run.

\begin{cousinweb}
\emph{discesa} is built on \emph{discendere} 'to go down', which is inherited from Latin \emph{dēscendere} (EtymDB), from \emph{dē} 'down' and \emph{scandere} 'to climb' (Lewis \& Short).
\end{cousinweb}

\begin{grammarlens}[title={What you've built}]
One more word for telling, explaining and describing at B1.
\end{grammarlens}

\subsection*{Your turn}
\begin{itemize}
\raggedright
  \item \emph{Say it:} \emph{la discesa}
  \item \emph{Say it:} \emph{la discesa}, once more
  \item \emph{Say it:} a sentence of your own with \emph{la discesa}
  \item \emph{Recall:} say the Italian for the one-way street, then the Italian for the climb, then say \emph{la discesa} again
\end{itemize}

\begin{tcolorbox}[breakable,colback=teal!4,colframe=teal!35!black,title={Before you move on}]
What does la discesa mean? (\textquotedblleft{}The descent, the downhill slope\textquotedblright{}.) Say it once more.
\end{tcolorbox}
\section[costeggiare]{costeggiare — to run along, to skirt (a coast or a river)}

\label{lesson:IT-C442-costeggiare}



\begin{quote}
Before the new one: say the Italian for the climb, then the Italian for the descent.
\end{quote}

\subsection*{You'll want to know: costeggiare}
\textbf{costeggiare} — \textquotedblleft{}to run along, to skirt (a coast or a river)\textquotedblright{}.

\begin{quote}
La strada costeggia il lago.
\end{quote}

\begin{quote}
\emph{The road runs along the lake.}
\end{quote}

\emph{La costa} — the coast.

\begin{cousinweb}
\emph{costeggiare} is built on \emph{costa} 'coast', which is inherited from Latin \emph{costa} 'rib, side' (EtymDB).
\end{cousinweb}

\begin{grammarlens}[title={What you've built}]
One more word for telling, explaining and describing at B1.
\end{grammarlens}

\subsection*{Your turn}
\begin{itemize}
\raggedright
  \item \emph{Say it:} \emph{costeggiare}
  \item \emph{Say it:} \emph{costeggiare}, once more
  \item \emph{Say it:} a sentence of your own with \emph{costeggiare}
  \item \emph{Recall:} say the Italian for the climb, then the Italian for the descent, then say \emph{costeggiare} again
\end{itemize}

\begin{tcolorbox}[breakable,colback=teal!4,colframe=teal!35!black,title={Before you move on}]
What does costeggiare mean? (\textquotedblleft{}To run along, to skirt (a coast or a river)\textquotedblright{}.) Say it once more.
\end{tcolorbox}
\section[sbucare]{sbucare — to come out, to pop out}

\label{lesson:IT-C442-sbucare}



\begin{quote}
Before the new one: say the Italian for the descent, then the Italian for to run along.
\end{quote}

\subsection*{You'll want to know: sbucare}
\textbf{sbucare} — \textquotedblleft{}to come out, to pop out\textquotedblright{}.

\begin{quote}
Dopo il vicolo siamo sbucati in piazza.
\end{quote}

\begin{quote}
\emph{After the alley we came out in the square.}
\end{quote}

\emph{Sbucare all'improvviso} — to pop up suddenly.

\begin{cousinweb}
\emph{sbucare} is built from \emph{s-} 'out of' and \emph{buca} 'hole': coming out of a hole. The origin of \emph{buca} is not settled.
\end{cousinweb}

\begin{grammarlens}[title={What you've built}]
One more word for telling, explaining and describing at B1.
\end{grammarlens}

\subsection*{Your turn}
\begin{itemize}
\raggedright
  \item \emph{Say it:} \emph{sbucare}
  \item \emph{Say it:} \emph{sbucare}, once more
  \item \emph{Say it:} a sentence of your own with \emph{sbucare}
  \item \emph{Recall:} say the Italian for the descent, then the Italian for to run along, then say \emph{sbucare} again
\end{itemize}

\begin{tcolorbox}[breakable,colback=teal!4,colframe=teal!35!black,title={Before you move on}]
What does sbucare mean? (\textquotedblleft{}To come out, to pop out\textquotedblright{}.) Say it once more.
\end{tcolorbox}
\section[aggirare]{aggirare — to go round, to get round}

\label{lesson:IT-C442-aggirare}



\begin{quote}
Before the new one: say the Italian for to run along, then the Italian for to come out.
\end{quote}

\subsection*{You'll want to know: aggirare}
\textbf{aggirare} — \textquotedblleft{}to go round, to get round\textquotedblright{}.

\begin{quote}
Aggiriamo il centro per evitare il traffico.
\end{quote}

\begin{quote}
\emph{We go round the centre to avoid the traffic.}
\end{quote}

Also to get round a problem or a rule.

\begin{cousinweb}
\emph{aggirare} is built from \emph{a-} and \emph{girare} 'to turn', which is inherited from Latin \emph{gyrāre}, from Greek \emph{gŷros} 'circle' (Lewis \& Short).
\end{cousinweb}

\begin{grammarlens}[title={What you've built}]
One more word for telling, explaining and describing at B1.
\end{grammarlens}

\subsection*{Your turn}
\begin{itemize}
\raggedright
  \item \emph{Say it:} \emph{aggirare}
  \item \emph{Say it:} \emph{aggirare}, once more
  \item \emph{Say it:} a sentence of your own with \emph{aggirare}
  \item \emph{Recall:} say the Italian for to run along, then the Italian for to come out, then say \emph{aggirare} again
\end{itemize}

\begin{tcolorbox}[breakable,colback=teal!4,colframe=teal!35!black,title={Before you move on}]
What does aggirare mean? (\textquotedblleft{}To go round, to get round\textquotedblright{}.) Say it once more.
\end{tcolorbox}
